%% ============================================================
%%  Appendix B — Chapter 3 (Bridge) Supporting Information
%%
%%  \include'd from thesis.tex under \appendix, after the
%%  Chapter 2 SI. Holds the full statement of the four
%%  methodological controls behind the degradation analysis;
%%  the bridge chapter keeps the result and a short summary.
%% ============================================================

\chapter{Supporting Information for Chapter~\ref*{ch:bridge}}\label{app:ch3_si}

This appendix records the methodological backbone of the degradation analysis of
\cref{sec:bridge_degradation}. The bridge chapter retains the result --- the
window collapse and the ohmic drift --- and the figures; the four controls that
make that result defensible against the confounds peculiar to this dataset are
stated in full here. The analysis is reproduced by
\texttt{scripts/\allowbreak bridge\_\allowbreak hybrane\_\allowbreak peo\_\allowbreak reproducibility.py}.

\section{The Four Methodological Controls}\label{app:ch3_controls}

The degradation analysis is built on a methodological backbone of four controls;
each was verified to change the conclusion if dropped.
\begin{enumerate}
  \item \textbf{Per-device weighting (mandatory).} The number of pixels measured per device collapses from sixteen to two over the campaign and is strongly anti-correlated with date (Spearman $\rho=-0.80$, $p\approx1\times10^{-15}$, $n=65$; \cref{fig:bridge_mechanism}a inset). Pixel-pooled statistics are therefore biased toward the early, densely measured devices. Every feature is reduced to \emph{one value per device} --- the median over that device's freshest-day curves --- before any test.
  \item \textbf{Standard-protocol corpus only} ($n\approx69$ devices, span 2021-02 to 2022-05): \SY{}/\Hybrane{}/\LiTf{} on silver, annealed at \SI{75}{\celsius} with no second stage, at the lead mass ratios (\Hybrane{}~\num{0.3}, \LiTf{}~\num{0.09}). This excludes the invasive experiments --- above all the high-temperature anneal batches that partially recovered devices a year in --- which would otherwise contaminate the feature-versus-date story.
  \item \textbf{Sweep-amplitude matching, with each feature read where it is expressed.} A $0\!\to\!+X\!\to\!0$ loop excites the device more at larger $X$, so it conducts more even when read back at the same voltage; amplitude is confounded with both date and material (\cref{fig:bridge_mechanism}b). Decisively, the \emph{window} features (on--off ratio, normalized area) only open up at high amplitude --- at $\sim$\SI{1.2}{\volt} they sit near unity for every device, early or late, and \emph{cannot} show a collapse --- so they are read in the $\sim$\SI{3}{\volt} stratum, while the \emph{conductivity} features are read in a tightly matched $\sim$\SI{1.2}{\volt} band.
  \item \textbf{Pre/post-inflection contrast and recovery-tail sensitivity.} Because the notes locate the behavioural inflection at NM\_v026 (2021-04-22), the window collapse is tested as a step --- early ($\leq$Apr 2021) versus later (Mann--Whitney) --- rather than as a whole-campaign correlation. Every conductivity trend is additionally reported with and without the deliberately-reverted recovery batches ($\geq$2022-03).
\end{enumerate}
The quantitative outcome of applying these controls --- the window collapse at
$\sim$\SI{3}{\volt} and the matched-amplitude ohmic drift --- is reported in
\cref{tab:bridge_degradation} and discussed in \cref{sec:bridge_degradation}.
